\documentclass[12pt]{article}
\usepackage[utf8]{inputenc}
\usepackage{amsmath}
\usepackage{amssymb}
\usepackage{amsthm}
\usepackage{geometry}
\geometry{margin=1in}

\newtheorem{theorem}{Theorem}
\newtheorem lemma}[theorem]{Lemma}
\newtheorem{corollary}[theorem]{Corollary}
\newtheorem{definition}{Definition}
\newtheorem{prop}[theorem]{Proposition}
\newtheorem{example}{Example}

\title{Mathematical Proofs: Q-Proof Quasicrystalline Blockchain Consensus}
\author{Fabian Leo Naressi\\
ORCID: 0009-0005-8045-0259\\
senemosiapuntozero@gmail.com}

\begin{document}
\maketitle

\section{Golden Ratio Properties}

\begin{definition}
The golden ratio $\phi$ is defined as the unique positive solution to the equation:
\begin{equation}
\phi^2 = \phi + 1
\label{eq:phi_def}
\end{equation}
yielding $\phi = \frac{1+\sqrt{5}}{2} \approx 1.618$ and its inverse $\phi^{-1} = \phi - 1 = \frac{\sqrt{5}-1}{2} \approx 0.618$.
\end{definition}

\begin{prop}
The golden ratio satisfies the identity:
\begin{equation}
\phi^n = F_n \phi + F_{n-1}
\label{eq:phi_fib}
\end{equation}
for all $n \in \mathbb{N}$, where $F_n$ denotes the $n$-th Fibonacci number.
\end{prop}

\begin{proof}
We prove by induction on $n$.

\textbf{Base case:} For $n=1$, $\phi^1 = F_1\phi + F_0 = 1\cdot\phi + 0 = \phi$. For $n=2$, $\phi^2 = \phi + 1 = F_2\phi + F_1 = 1\cdot\phi + 1$.

\textbf{Inductive step:} Assume $\phi^k = F_k\phi + F_{k-1}$ and $\phi^{k-1} = F_{k-1}\phi + F_{k-2}$. Then:
\begin{align}
\phi^{k+1} &= \phi \cdot \phi^k = \phi(F_k\phi + F_{k-1}) \\
&= F_k\phi^2 + F_{k-1}\phi = F_k(\phi + 1) + F_{k-1}\phi \\
&= (F_k + F_{k-1})\phi + F_k = F_{k+1}\phi + F_k
\end{align}
which completes the induction. $\square$
\end{proof}

\begin{corollary}
The ratio of consecutive Fibonacci numbers converges to $\phi$:
\begin{equation}
\lim_{n\to\infty} \frac{F_{n+1}}{F_n} = \phi
\label{eq:fib_converge}
\end{equation}
\end{corollary}

\begin{prop}
The golden ratio conjugate $\hat{\phi} = \frac{1-\sqrt{5}}{2}$ satisfies $|\hat{\phi}| < 1$, enabling efficient computation of Fibonacci numbers via Binet's formula:
\begin{equation}
F_n = \frac{\phi^n - \hat{\phi}^n}{\sqrt{5}}
\label{eq:binet}
\end{equation}
\end{prop}

\section{$\phi$-Spacing Topology}

\begin{definition}
A $\phi$-spacing sequence $\{d_i\}$ in a distributed network is defined recursively as:
\begin{equation}
d_i = d_0 \cdot \phi^i, \quad i \in \mathbb{Z}
\label{eq:phi_spacing}
\end{equation}
where $d_0 > 0$ is the base spacing parameter.
\end{definition}

\begin{prop}
The $\phi$-spacing sequence is \emph{deterministically aperiodic}: there exists no non-zero $T \in \mathbb{R}$ such that $d_{i+T} = d_i$ for all $i$.
\end{prop}

\begin{proof}
Assume for contradiction that a period $T > 0$ exists. Then for all $i$:
\begin{equation}
d_0 \phi^{i+T} = d_0 \phi^i \implies \phi^T = 1
\end{equation}
But $\phi > 1$ is irrational, so $\phi^T = 1$ implies $T = 0$, contradiction. $\square$
\end{proof}

\begin{prop}
The spacing ratios between consecutive nodes are irrational:
\begin{equation}
\frac{d_{i+1}}{d_i} = \phi \notin \mathbb{Q}
\label{eq:irrational_ratio}
\end{equation}
\end{prop}

\begin{corollary}
For any rational threshold $q = a/b \in \mathbb{Q}$, the equation $\phi^i \equiv q \pmod{1}$ has at most one solution, ensuring unpredictability of node positions modulo any rational base.
\end{corollary}

\section{Consensus Security Proofs}

\begin{prop}[51\% Attack Resistance]
Let $S \subset V$ be a coalition attempting a 51\% attack, where $|S|/|V| \geq 0.5$. In a $\phi$-spaced topology, the probability that $S$ forms a connected adversarial subgraph is bounded by:
\begin{equation}
P(\text{attack}) \leq \prod_{i=1}^{|S|} \frac{1}{\phi^i} = \phi^{-|S|(|S|+1)/2}
\label{eq:attack_bound}
\end{equation}
\end{prop}

\begin{proof}
For node $i$ to connect to the adversarial coalition, its $\phi$-spacing position must align with members of $S$. The alignment probability decreases geometrically with $\phi$ due to the irrational spacing. The adversarial coalition can only form if each node's position satisfies irrational constraints that match $S$. By the properties of $\phi$-spacing, the joint probability factors as the product of individual alignment probabilities. $\square$
\end{proof}

\begin{corollary}
For a network of $n$ validators with $|S| = 0.51n$, the attack probability scales as:
\begin{equation}
P(\text{attack}) \leq \phi^{-0.1301n^2} = O(e^{-cn^2})
\label{eq:exp_bound}
\end{equation}
for some constant $c > 0$, providing exponential security.
\end{corollary}

\section{Zero-Point Energy Formalism}

\begin{prop}
The total energy of the quasicrystalline consensus network is given by:
\begin{equation}
E_{\text{total}} = E_{\text{interaction}} + E_{\text{ZPE}}
\label{eq:energy_total}
\end{equation}
where the interaction energy is quasiperiodic and the zero-point energy provides ground-state stability.
\end{prop}

\begin{prop}
The ground state energy (zero-point energy) of the system is:
\begin{equation}
E_{\text{ZPE}} = \frac{1}{2}\sum_{i} \hbar\omega_i = \frac{\hbar}{2}\sum_{i} k_i v_F
\label{eq:zpe}
\end{equation}
where $\omega_i$ are the characteristic frequencies determined by the $\phi$-spacing.
\end{prop}

\begin{prop}
The fluctuation-dissipation relation for ZPE-guided consensus is:
\begin{equation}
\chi(\omega) = \frac{1}{\hbar\omega + i\epsilon}
\label{eq:fluctuation}
\end{equation}
ensuring the system remains in a metastable ground state configuration.
\end{prop}

\section{Self-Healing Properties}

\begin{prop}
Let $G = (V, E)$ be the $\phi$-spaced validator graph. After node failure, the system achieves equilibrium within $O(\log |V|)$ steps through energy minimization.
\end{prop}

\begin{prop}
The convergence time satisfies:
\begin{equation}
T_{\text{heal}} = T_0 \cdot \log_{\phi}(N)
\label{eq:heal_time}
\end{equation}
where $N$ is the number of validators and $T_0$ is the base healing time.
\end{prop}

\section{Sybil Resistance}

\begin{prop}
In the $\phi$-spacing topology, a Sybil attack requires modifying the irrational spacing by at least $\epsilon > 0$. The detection probability is:
\begin{equation}
P_{\text{detect}} = 1 - \prod_{i=1}^{k} \left|\hat{\phi}^i - \hat{\phi}_{\text{new}}^i\right|
\label{eq:sybil_detect}
\end{equation}
for $k$ observed connections.
\end{prop}

\section{Scalability Analysis}

\begin{prop}
The propagation latency in a $\phi$-spaced network of $N$ validators scales as:
\begin{equation}
T_{\text{prop}}(N) = T_0 \cdot \log_{\phi}(N) = O(\log N)
\label{eq:latency}
\end{equation}
compared to $T_{\text{classic}}(N) = O(N)$ for traditional periodic topologies.
\end{prop}

\begin{corollary}
For $N = 2^{10} \approx 1024$ validators, the speedup factor is:
\begin{equation}
\frac{T_{\text{classic}}}{T_{\phi}} = \frac{N}{\log_{\phi}(N)} \approx \frac{1024}{13.2} \approx 77.5\times
\label{eq:speedup}
\end{equation}
\end{corollary}

\section{Quantum Resistance}

\begin{prop}
The ML-DSA signature scheme provides security against quantum adversaries under the Module-LWE assumption. The success probability of a quantum forger is bounded by:
\begin{equation}
\text{Succ}_{\text{ML-DSA}}^{\text{quantum}} \leq O(q^{-1/2})
\label{eq:ml_dsa}
\end{equation}
where $q$ is the lattice modulus.
\end{prop}

\begin{thebibliography}{99}

\end{thebibliography}

\end{document}
